\section{古典的 Poisson 双有理写像の blow-up 表示}
\label{sec:classical-factorizations}

標数零の滑らかな完備多様体間の双有理写像には、滑らかな中心の
blow-up と blow-down による弱分解がある
\cite[定理0.1.1]{AKMW}。
Poisson 構造を各中間模型上で正則に保つには、各中心に追加条件が必要になる。
本節では使用する中心を具体的に指定し、共通の滑らかな Poisson 模型を構成する。

\subsection{曲線中心と点中心の持ち上げ条件}

\begin{proposition}\label{prop:poisson-centres}
滑らかな三次元 Poisson 多様体 $(X,\pi)$ と滑らかな中心 $Z$ を考える。
\begin{enumerate}
\item $Z$ が曲線ならば、$\pi$ が $\operatorname{Bl}_Z X$ に正則に持ち上がる
必要十分条件は $\pi|_Z=0$ である。
\item $Z=\{p\}$ が点ならば、必要十分条件は $\pi(p)=0$ と、$p$ を原点とする
局所座標での一次項が
\[
 \pi_1=\bigl(x\partial_x+y\partial_y+z\partial_z\bigr)\wedge V
\]
と書けることである。ここで $V$ は定係数ベクトル場である。
\end{enumerate}
いずれの場合も持ち上げは一意的で、blow-down 射は Poisson 射となる。
\end{proposition}

\begin{proof}
曲線の場合、$Z=\{v=w=0\}$ として
\[
 \pi=A\partial_v\wedge\partial_w+
 B\partial_w\wedge\partial_x+C\partial_x\wedge\partial_v
\]
と書く。$w=v\tau$ のチャートでは
\[
 \widetilde\pi=\frac{A}{v}\partial_v\wedge\partial_\tau+
 C\partial_x\wedge\partial_v-
 \frac{B+\tau C}{v}\partial_x\wedge\partial_\tau.
\]
例外因子上の任意の $\tau$ で正則となる条件は
$A|_Z=B|_Z=C|_Z=0$ である。他方のチャートも同じ条件を与える。

点の場合、$A=\{y,z\}$、$B=\{z,x\}$、$C=\{x,y\}$ と置く。
$u=y/x$、$v=z/x$ のチャートでは
\[
 \{x,u\}=C/x,\qquad \{x,v\}=-B/x,\qquad
 \{u,v\}=\frac{xA+yB+zC}{x^3}.
\]
三つのチャートの正則性から定数項は全て零となる。
$A_1,B_1,C_1$ を一次項とすると、残る条件は
\[
 xA_1+yB_1+zC_1=0.
\]
一次式のこの関係は、$(A_1,B_1,C_1)$ の係数行列が歪対称であることと同値であり、
ちょうど命題の $\pi_1$ を与える。二次以上の項については、上の最後の分子は
三次以上なので、各 blow-up チャートへの代入後に分母が消える。
これで十分性も従う。
一意性と Jacobi 恒等式は中心の補集合の稠密性による。
\end{proof}

この条件は Polishchuk \cite[定理8.2、命題8.3]{Polishchuk1997} の
退化した余法 Lie 代数という条件の三次元での表示である。
例えば $\pi=x\partial_y\wedge\partial_z$ は原点で零であるが、
$\{y/x,z/x\}=1/x$ なので原点の blow-up には持ち上がらない。
一方、以下の二次曲面 pencil の点中心では一次項の条件も満たされる。

\subsection{二次曲面 pencil と垂直型の Poisson 双有理同値}
\label{sec:birational-rat22-vertical}

この節では基礎体を $\mathbf C$ とし、双有理写像が誘導する関数体の同型が
Poisson 括弧を保存するとき、二つの Poisson 多様体は Poisson 双有理同値であるという。
ここで用いる双有理同値は、滑らかな共通モデルからの正則 Poisson 射によって実現される。

\begin{theorem}
\label{thm:rat22-vertical}
$\mathbf P^3$ 上の二次式 $q_0,q_1$ の共通零点
$E=\{q_0=q_1=0\}$ が滑らかな完全交叉曲線であるとし、
\[
 \alpha=q_0\,dq_1-q_1\,dq_0
 \in H^0\bigl(\mathbf P^3,\Omega^1_{\mathbf P^3}(4)\bigr)
\]
に対応する Poisson 双ベクトルを $\pi$ とする。
任意の $p\in E$ に対し、滑らかな平面三次曲線
$C\subset\mathbf P^2$ と、滑らかな射影三次元多様体 $Z$ 上の正則 Poisson 構造
$\widetilde\pi$ が存在し、
\[
 (\mathbf P^3,\pi)
 \xleftarrow{\ \beta\ }
 (Z,\widetilde\pi)
 \xrightarrow{\ \psi\ }
 (\mathbf P^1\times\mathbf P^2,\pi_C)
\]
の両射は双有理かつ正則な Poisson 射となる。
ここで $\pi_C$ は $\mathbf P^2$ 上の三次式 $C$ に対応する Poisson 構造の
$\mathbf P^1\times\mathbf P^2$ への垂直な延長である。
また $C$ は $E$ と同型であり、具体的に
\[
 Z=\operatorname{Bl}_{E'}\bigl(\operatorname{Bl}_p\mathbf P^3\bigr)
 \cong\operatorname{Bl}_{\Gamma}(\mathbf P^1\times\mathbf P^2),
\]
と取れる。ここで $E'$ は $E$ の狭義変換であり、
$\Gamma$ は $(1,1)$ 型および $(1,2)$ 型因子の滑らかな完全交叉である。
第二射影は $\Gamma\cong C$ を与え、第一射影 $\Gamma\to\mathbf P^1$ の次数は $2$ である。
逆に、任意の非零三次式で定義される滑らかな平面三次曲線と、その垂直型 Poisson 構造は、
この構成から得られる。
\end{theorem}

\begin{proof}
$p=[1:0:0:0]$ とする斉次座標 $[x_0:x_1:x_2:x_3]$ を取り、
$y=[x_1:x_2:x_3]$ と書く。
$q_i(p)=0$ より、一意的に
\[
 q_i=x_0\ell_i(y)+r_i(y)\qquad(i=0,1)
\]
と書ける。ここで $\ell_i$ は一次式、$r_i$ は二次式である。
$E$ の $p$ における滑らかさから $\ell_0,\ell_1$ は一次独立である。
\[
 C(y)=\ell_1(y)r_0(y)-\ell_0(y)r_1(y)
\]
と置き、次の双有理写像を考える。
\[
 \varphi\colon\mathbf P^3\dashrightarrow\mathbf P^1\times\mathbf P^2,
 \qquad x\longmapsto\bigl([q_0(x):q_1(x)],y\bigr).
 \tag{\ref{sec:birational-rat22-vertical}.1}
\]
逆写像を斉次座標で表すため、$\mathbf P^1$ の座標を $[a:b]$ とし、
\[
 D=a\ell_1-b\ell_0,\qquad N=br_0-ar_1
\]
と置く。このとき逆写像は
\[
 \chi\colon\mathbf P^1\times\mathbf P^2\dashrightarrow\mathbf P^3,
 \qquad ([a:b],[y_1:y_2:y_3])
 \longmapsto[N:Dy_1:Dy_2:Dy_3]
\]
である。実際、$q_0(N,Dy)=aDC$, $q_1(N,Dy)=bDC$ が成り立つ。
$a\ne0$ の開集合で $u=b/a$ とすれば、これは
$x_0=(u r_0-r_1)/(\ell_1-u\ell_0)$ と同値である。

まず Poisson 双ベクトルの変換を計算する。
$x_3=1$ とし、$v=x_1/x_3$, $w=x_2/x_3$ を使う。
このアフィン座標で $\iota_\pi(dx_0\wedge dv\wedge dw)=\alpha$ とすると、
\[
 \alpha=q_0^2\,du,
 \qquad
 \frac{\partial u}{\partial x_0}
 =\frac{\ell_1q_0-\ell_0q_1}{q_0^2}
 =\frac{C(v,w,1)}{q_0^2}.
\]
$u$ は Casimir であり、$\{v,w\}_\pi$ は $\alpha$ の $dx_0$ 係数だから、
\[
 \varphi_*\pi
 =C(v,w,1)\,\partial_v\wedge\partial_w.
 \tag{\ref{sec:birational-rat22-vertical}.2}
\]
これは $\mathbf P^1\times\mathbf P^2$ 全体で正則である。
この等式は双ベクトルそのものの変換公式である。

次に $C=0$ の滑らかさを確認する。
$E$ は楕円曲線であり、$H=\mathcal O_E(1)$ は次数 $4$ の直線束である。
$p$ からの射影が $E$ 上に誘導する写像は完全線形系 $|H(-p)|$ で与えられる。
次数 $3$ の直線束は楕円曲線上で非常に豊富なので、これは $E$ を滑らかな平面三次曲線に埋め込む。
$q_0=q_1=0$ から $x_0$ を消去すると像の方程式は $C=0$ となる。
したがって、この三次式は滑らかな曲線を定義する。

$B=\operatorname{Bl}_p\mathbf P^3$ と置く。
$p$ からの射影は $B\to\mathbf P^2$ に延長する。
$p$ の近傍で $s=q_0$, $t=q_1$ を適切な局所自明化で表し、
$E=\{s=t=0\}$ となるエタール局所座標 $(s,t,e)$ を取る。
$\pi$ は、正則可逆関数 $h$ を用いて
\[
 \pi=h(s,t,e)\,(s\partial_s+t\partial_t)\wedge\partial_e
 \tag{\ref{sec:birational-rat22-vertical}.3}
\]
と書ける。符号は $h$ に含めている。
点の blowup の $s$ チャート
$(s,t,e)=(a,ab,ac)$ では
\[
 (s\partial_s+t\partial_t)\wedge\partial_e
 =\partial_a\wedge\partial_c,
\]
$e$ チャート $(s,t,e)=(ea,eb,e)$ では
\[
 (s\partial_s+t\partial_t)\wedge\partial_e
 =(a\partial_a+b\partial_b)\wedge\partial_e.
\]
$t$ チャートも同様なので、$\pi$ は $B$ 上の正則 Poisson 構造に持ち上がる。
$E'$ は $e$ チャートの $a=b=0$ であり、持ち上げた双ベクトルは $E'$ 上で消える。
さらに $E'$ を blowup し、例えば $b=ac$ と置くと
\[
 (a\partial_a+b\partial_b)\wedge\partial_e
 =a\partial_a\wedge\partial_e.
\]
よって $Z=\operatorname{Bl}_{E'}B$ にも正則に持ち上がる。
各段階の Jacobi 恒等式は稠密開集合で成立するため全体で成立する。

$B$ 上の pencil の基底イデアルは、点の blowup の例外因子を共通因子として除くと
$E'$ のイデアルになる。実際、$s$ チャートでは残余イデアルは $(1,b)$ であり、
$e$ チャートでは $(a,b)$ である。
したがって $Z$ 上で $[q_0:q_1]$ も射に延長し、
射影 $Z\to\mathbf P^2$ と合わせて双有理射
$\psi\colon Z\to\mathbf P^1\times\mathbf P^2$ を得る。
式~(\ref{sec:birational-rat22-vertical}.2) より、$\psi$ は稠密開集合で Poisson 射である。
双方の構造と射が正則であるから、$\psi$ は全体で Poisson 射となる。
$\beta$ についても同様である。

$\psi$ を blowup として同定する。
$\chi$ の四つの座標は $\mathcal O(1,2)$ の切断であり、その基底イデアルは
\[
 (N,Dy_1,Dy_2,Dy_3)=(N,D)
\]
である。ここで右辺との等式はイデアル層の等式であり、各点で少なくとも一つの $y_j$ が
可逆であることによる。
したがって $\chi$ のグラフの閉包は
$\operatorname{Bl}_{\Gamma}(\mathbf P^1\times\mathbf P^2)$ となる。ただし
\[
 \Gamma=\{D=N=0\}.
\]
一方、$B$ は $p$ からの射影のグラフの閉包であり、
$Z=\operatorname{Bl}_{E'}B$ はその上の残余 pencil のグラフの閉包である。
従って $Z$ は $\varphi$ のグラフの閉包そのものである。
互いに逆な双有理写像のグラフの閉包は因子の順序を交換すると一致するため、
\[
 Z\cong\operatorname{Bl}_{\Gamma}(\mathbf P^1\times\mathbf P^2)
\]
を得る。この同型のもとで $\psi$ は blowup の射である。

$\Gamma$ の滑らかさも直接確認できる。
$p_0=\{\ell_0=\ell_1=0\}$ と置くと、$C(p_0)=0$ である。
$y\ne p_0$ では $D=0$ が $[a:b]=[\ell_0(y):\ell_1(y)]$ を定め、
$N=0$ は $C(y)=0$ と同値になる。
$C$ の滑らかさから $(r_0(p_0),r_1(p_0))\ne(0,0)$ である。
例えば $r_0(p_0)\ne0$ なら、$a=1$ とするチャートで
$N=0$ から $b=r_1/r_0$ を得て、残りの方程式は
$D=C/r_0=0$ となる。
従って第二射影はスキームの同型 $\Gamma\cong C$ を与える。
特に $\Gamma$ は滑らかな楕円曲線であり、表示通り $(1,1),(1,2)$ 型の完全交叉である。
第一射影は $C$ 上の線形系 $|\mathcal O_C(1)(-p_0)|$ で与えられる。
この直線束の次数は $2$ で、完全線形系は基点を持たないので、
$\Gamma\to\mathbf P^1$ の次数は $2$ である。

逆に、滑らかな平面三次式 $C(y)$ と点 $p_0\in\{C=0\}$ を取る。
$p_0$ のイデアルを生成する独立な一次式 $\ell_0,\ell_1$ を選ぶと、
$C\in(\ell_0,\ell_1)$ より二次式 $r_0,r_1$ を用いて
\[
 C=\ell_1r_0-\ell_0r_1
\]
と書ける。$q_i=x_0\ell_i+r_i$ と置く。
$p=[1:0:0:0]$ では二つの微分 $dq_0,dq_1$ が一次独立である。
$y\ne p_0$ では $q_0=q_1=0$ は $C(y)=0$ と $x_0$ を一意に定める一次方程式に同値であり、
この部分は $C\setminus\{p_0\}$ と同型である。
$C$ の $p_0$ における滑らかさから
$(r_0(p_0),r_1(p_0))\ne(0,0)$ であり、$y=p_0$ の上に $p$ 以外の点はない。
以上により $\{q_0=q_1=0\}$ は滑らかな完全交叉曲線となる。
先の変換公式は、指定した三次式 $C$ の係数を含めてそのまま与える。
\end{proof}

\begin{remark}
二次曲面 pencil 成分 $R(2,2)$ は
$\operatorname{Gr}(2,H^0(\mathbf P^3,\mathcal O(2)))$ に対応し、
$\dim H^0(\mathbf P^3,\mathcal O(2))=10$ なので、アフィン Poisson 空間内での次元は
$2(10-2)+1=17$ である。
垂直成分 $V_{\mathrm{vert}}$ は
$H^0(\mathbf P^2,\mathcal O(3))$ であり、アフィン次元は $10$ である。
定理~\ref{thm:rat22-vertical} は両成分の滑らかな楕円曲線を伴う開集合の間の
Poisson 双有理同値を構成する。写像の選択には、基底楕円曲線上の点の選択が含まれる。
\end{remark}

\subsection{楕円曲線に沿う基本変換と三重積のペンシル}
$S=\mathbb P^1_y\times\mathbb P^1_z$, $X=\mathbb P^1_x\times S$ とおく。
$s,t\in H^0(X,\mathcal O(1,1,1))$ を
\[
 s=xa(y,z)+b(y,z),\qquad t=xc(y,z)+d(y,z),\qquad
 a,b,c,d\in H^0(S,\mathcal O(1,1))
\]
と書き、$\alpha=s\,dt-t\,ds$ に対応する Poisson 構造を $\pi$ とする。
\[
 D=bc-ad\in H^0(S,\mathcal O(2,2)),\qquad u=t/s
\]
とおけば、関数体上で $u$ は Casimir であり
\[
 \{u,y\}=\{u,z\}=0,\qquad
 \{y,z\}=s\partial_xt-t\partial_xs=D(y,z).
\]
従って有理写像
\[
 \Phi\colon X\dashrightarrow X'=\mathbb P^1_u\times S,
 \qquad (x,y,z)\longmapsto(t/s,y,z)
\]
は $\pi$ を実際の双ベクトル $D(y,z)\partial_y\wedge\partial_z$ に送る。
変換式は双ベクトル自体の等式である。
$D\ne0$ なら逆写像は
\[
 x=\frac{d-ub}{ua-c}
\]
なので $\Phi$ は双有理である。

一般の $s,t$ に対して $E=\{s=t=0\}\subset X$ は滑らかな楕円曲線である。
また $\Gamma=\{D=0\}\subset S$ は滑らかな $(2,2)$ 曲線である。
$a,b,c,d$ は同時には零にならず、$\Gamma$ 上で行列
$\left(\begin{smallmatrix}a&b\\c&d\end{smallmatrix}\right)$ は階数一なので、
射影 $E\to\Gamma$ は同型になる。
$[x_0:x_1]$, $[u_0:u_1]$ を $x=x_1/x_0$, $u=u_1/u_0$ に対応する座標とすれば、
$\Phi$ の閉グラフは
\[
 B=\{x_1(u_0c-u_1a)+x_0(u_0d-u_1b)=0\}
 \subset\mathbb P^1_x\times\mathbb P^1_u\times S.
\]
これは同時に
\[
 B\cong\operatorname{Bl}_{E}X
 \cong\operatorname{Bl}_{E'}X',\qquad
 E'=\{u_0c-u_1a=u_0d-u_1b=0\}
\]
である。実際、どちらの射影に関しても二つの切断の比のグラフであり、
それらの共通零点は滑らかな余次元二の完全交叉である。
同様に $E'\to\Gamma$ は同型で、$B$ は滑らかである。

命題\ref{prop:poisson-centres}の曲線中心の条件を適用する。

$E$ 上では $s=t=0$ より $\pi=0$ であり、$E'$ 上では $D=0$ である。
従って $B$ 上には正則 Poisson 構造があり、上の二つの射影はいずれも
Poisson 射である。二つの持ち上げは稠密開集合で等しいので一致する。
一般の $u$ に対する葉の閉包は、元の側では $(1,1,1)$ 型の
次数六の del Pezzo 曲面、変換先では $\mathbb P^1\times\mathbb P^1$ である。
$E'\to\mathbb P^1_u$ の次数は二であり、共通モデルの上では後者を二点で
ブローアップして前者を得る。

この構成は一般の $(2,2)$ 型の Poisson 双ベクトルを与える。
実際、$V=H^0(S,\mathcal O(1,1))$, $W=H^0(S,\mathcal O(2,2))$ とおくと、
行列式写像 $V^4\to W$, $(a,b,c,d)\mapsto bc-ad$ の微分は
\[
 (\dot a,\dot b,\dot c,\dot d)\longmapsto
 c\dot b+b\dot c-d\dot a-a\dot d.
\]
$a,b,c,d$ が $V$ を張る点では、その像は $V\cdot V=W$ である。
従って行列式写像は優越的であり、$D$ が滑らかである開条件と
$a,b,c,d$ が同時には消えない開条件を満たす一般点で、上の構成が適用できる。

\subsection{平面と二重積の曲面リンク}
\label{subsec:classical-surface-link}
相異なる $p,q\in\mathbb P^2$ と、$r\in\mathbb P^1\times\mathbb P^1$ に対して
標準的な曲面の同型
\[
 S_0=\operatorname{Bl}_{p,q}\mathbb P^2
 \cong\operatorname{Bl}_{r}(\mathbb P^1\times\mathbb P^1)
\]
がある。左側では $p,q$ を結ぶ直線の狭義変換を右側の $r$ へ収縮する。
座標を
\[
 p=[0:1:0],\quad q=[0:0:1],\quad
 [T:Y:Z]=[1:y:z],\quad r=(\infty,\infty)
\]
と取れば、有理写像は共通アフィン平面上で恒等写像である。
このリンクに $\mathbb P^1_x$ を掛けた三次元のリンクを考える。

三重積上の双ベクトルを
\[
 \rho=A(y,z)\partial_y\wedge\partial_z+
       B(x,z)\partial_z\wedge\partial_x+
       C(x,y)\partial_x\wedge\partial_y
\]
と書く。$A=\sum_{0\le i,j\le2}a_{ij}y^iz^j$ とし、
$B=b_0(x)+b_1(x)z+b_2(x)z^2$、
$C=c_0(x)+c_1(x)y+c_2(x)y^2$ とする。
この同じ有理双ベクトルが $\mathbb P^1\times\mathbb P^2$ に正則に延びる条件は
\begin{equation}\label{eq:surface-link-conditions}
 a_{22}=0,\qquad b_2(x)\equiv0,\qquad c_2(x)\equiv0.
\end{equation}
これは $\rho|_{\mathbb P^1\times\{r\}}=0$ と同値である。
直接には、$a_{22}=0$ は $A$ の全次数が三以下となる条件であり、
残り二条件は混合ベクトル場が平面の射影ベクトル場となる条件である。
一般の射影ベクトル場の二次部分は
$\ell(y,z)(y\partial_y+z\partial_z)$ なので、
各係数が一方の変数だけに依存する本表示では二次部分は消えなければならない。

逆向きには、$\mathbb P^1_x\times\mathbb P^2$ 上の双ベクトルを
\[
 \rho=F\,\partial_y\wedge\partial_z+
       \sum_{i=0}^2 x^i\partial_x\wedge V_i
\]
と書く。ここで $F$ は斉次平面三次式のアフィン表示であり、$V_i$ は
平面の射影ベクトル場である。三重積への正則延長の必要十分条件は
\[
 F(p)=F(q)=0,\qquad V_i(p)=V_i(q)=0\quad(0\le i\le2).
\]
すなわち $\mathbb P^1\times\{p,q\}$ 上で双ベクトルが零になる条件である。
固定した座標の下で両側の条件を満たす双ベクトルの空間はともに二十次元で、
共通アフィンチャートでの恒等写像がそれらを線形に同一視する。
Jacobi 方程式も同一視される。
これらの条件が成り立つとき $\mathbb P^1\times S_0$ の上への持ち上げは正則で、
両射影が Poisson 射になる。

特に $F=0$ が滑らかな平面三次曲線 $\Gamma_2$ で、混合項が零ならば、
$p,q\in\Gamma_2$ を相異なる一般の二点に選べる。
その狭義変換は $S_0$ 上でも滑らかな楕円曲線であり、
$p,q$ を結ぶ直線の狭義変換と一点で横断的に交わる。
従って収縮後の曲線 $\Gamma_{11}$ は滑らかな $(2,2)$ 曲線である。
クラスを用いると、$H$ を平面の直線、$E_p,E_q$ を例外曲線として
\[
 F_1=H-E_p,\quad F_2=H-E_q,\quad E_r=H-E_p-E_q,
 \qquad 3H-E_p-E_q=2F_1+2F_2-E_r
\]
となる。これが $(2,2)$ 型を与える。
逆に一般の滑らかな $(2,2)$ 曲線上の一般点を $r$ に選べば、
二本の座標線を収縮して滑らかな平面三次曲線を得る。
従って純粋な曲面 Poisson 構造については、平面三次曲線の一般点と
二重積の $(2,2)$ 曲線の一般点が、正則 Poisson 共通モデルを介して対応する。

上の曲面リンクを三重積の一般のペンシルに直接適用するには、
\eqref{eq:surface-link-conditions} を満たす特殊化が必要である。
実際、一般のペンシルでは三つの辺の $3\times3$ テンソルはいずれも階数三である。
例えば $s=1+xyz$, $t=x+y+z$ では各辺の係数は
$1-yz(y+z)$ とその置換であり、係数行列の行列式は非零である。
$\mathbb P^1_x\times\{(y_0,z_0)\}$ 上で混合項が恒等的に消える条件は、
階数三のテンソルを $y_0$ または $z_0$ における非零評価写像で縮約して零にする条件なので、
成立しない。ペンシルから純粋な曲面型へ先に楕円曲線に沿う基本変換を行うことで、
その後に曲面リンクを適用できる。

\subsection{三つの標準的なコンパクト化のトーラス型}
$T=(\mathbb C^*)^3$ の座標を $x_1,x_2,x_3$ とする。
$\lambda=(\lambda_{ij})$ を定数歪対称行列として
\[
 \pi_\lambda=\sum_{i<j}\lambda_{ij}
 (x_i\partial_{x_i})\wedge(x_j\partial_{x_j})
\]
と置く。三つのベクトル場は可換なので Jacobi 条件が成り立つ。
任意の滑らかなトーリックコンパクト化でこれらは大域正則ベクトル場に
延長する。従って $\pi_\lambda$ は $\mathbb P^3$、
$\mathbb P^1\times\mathbb P^2$、$(\mathbb P^1)^3$ のそれぞれで正則であり、
共通の稠密トーラス上で恒等写像によって同じ Poisson 双ベクトルを与える。

この対応には滑らかな共通 Poisson 解消が存在する。
$N=\mathbb Z^3$ において $v_i=e_i$ $(1\le i\le3)$、
$v_0=-e_1-e_2-e_3$ とする。非空真部分集合 $I\subset\{0,1,2,3\}$
について $v_I=\sum_{i\in I}v_i$ を取り、部分集合の鎖に対応する錐を
すべて取ると $A_3$ 型 permutohedral fan を得る。
$N_{\mathbb R}$ の実座標を $\xi_1,\xi_2,\xi_3$ とすると、
その極大錐は $0,\xi_1,\xi_2,\xi_3$ の大小関係を一つに固定する錐であり、
その各原始生成元は $N$ の基底をなす。この扇は
$\mathbb P^3$ の扇、座標直交象限からなる $(\mathbb P^1)^3$ の扇、
および、$\mathbb P^1$ の方向を $e_3$、$\mathbb P^2$ の三つの半直線を
$e_1,e_2,-e_1-e_2$ とする積の扇を細分する。
対応する滑らかな射影トーリック三次元多様体 $Z_{\rm per}$ から三つの
コンパクト化への固有双有理トーリック射が得られる。
$\pi_\lambda$ は $Z_{\rm per}$ 上でも正則であり、各射が稠密トーラス上で
Poisson であることから全体で Poisson である。
三つの射はいずれも滑らかなトーラス不変中心の blow-up の合成である。
その中心を以下に明示する。滑らかな扇の錐
$\tau=\langle r_1,\ldots,r_c\rangle$ に対応する軌道閉包を blow-up すると、
$r_\tau=r_1+\cdots+r_c$ を加える star subdivision になる。
三次元錐を中心にすれば点、二次元錐を中心にすれば曲線の blow-up である。
中心とそれに対応する錐は、各段階の模型上で取る。

\begin{proposition}\label{prop:toric-explicit-blowups}
共通模型 $Z_{\rm per}$ は次の三つの列のいずれによっても得られる。
\begin{enumerate}
\item $\mathbb P^3$ から四つの座標点を blow-up し、続いて六つの座標直線の
狭義変換を blow-up する。対応する中心錐は、まず
$\langle v_i,v_j,v_k\rangle$ $(i<j<k)$ の四つ、続いて
$\langle v_i,v_j\rangle$ $(i<j)$ の六つである。
\item $(\mathbb P^1)^3$ から、二つの点
$\langle e_1,e_2,e_3\rangle$、$\langle-e_1,-e_2,-e_3\rangle$
を blow-up する。続いて六つの曲線
$\langle e_i,e_j\rangle$ および $\langle-e_i,-e_j\rangle$ $(i<j)$
の狭義変換を blow-up する。
\item $\mathbb P^1\times\mathbb P^2$ の扇を上と同じ規約で取り、
$c=-e_1-e_2$、$w=-e_1-e_2-e_3$ と置く。
次の表の順序で、一点と八曲線を blow-up する。
\[
\begin{array}{c|c|c}
 \text{段階}&\text{中心錐}&\text{新しい半直線の原始生成元}\\\hline
 1&\langle e_1,e_2,e_3\rangle&e_1+e_2+e_3\\
 2&\langle c,-e_3\rangle&w\\
 3&\langle c,e_2\rangle&-e_1\\
 4&\langle w,e_2\rangle&-e_1-e_3\\
 5&\langle c,e_1\rangle&-e_2\\
 6&\langle w,e_1\rangle&-e_2-e_3\\
 7&\langle e_2,e_3\rangle&e_2+e_3\\
 8&\langle e_1,e_3\rangle&e_1+e_3\\
 9&\langle e_1,e_2\rangle&e_1+e_2
\end{array}
\]
\end{enumerate}
各列の全中間模型で $\pi_\lambda$ は正則であり、各 blow-down は Poisson 射である。
従って三つの端点の間のトーラス型の双有理写像は、この共通模型を通る
Poisson blow-up と blow-down の列として表せる。
\end{proposition}
\begin{proof}
各列について star subdivision の規則を適用する。
以下では $\sigma(1)$ を錐 $\sigma$ の半直線の原始生成元の集合とする。
中心錐 $\tau$ を含む最大錐 $\sigma$ は、$r\in\tau(1)$ ごとに
\[
 \bigl\langle r_\tau,\ \sigma(1)\setminus\{r\}\bigr\rangle
\]
で置き換え、それ以外の最大錐は保つ。
上の各中心は指定された段階で実際に錐をなし、新しく得られる最大錐の
三本の生成元の行列式は $\pm1$ である。
また各段階の二次元面の線形包は
$\xi_i=0$ または $\xi_i=\xi_j$ の形になる。
これらは、表および二つの対称な列の生成元について外積を取れば確認できる。
従って $0,\xi_1,\xi_2,\xi_3$ の大小関係による扇は、全中間扇を細分する。

三つの初期扇の最大錐の個数はそれぞれ $4,8,6$ である。
三次元錐の star subdivision は一つの最大錐を三つに置き換える。
二次元錐は滑らかな完備三次元扇のちょうど二つの最大錐に含まれ、
それぞれを二分する。従っていずれの操作も最大錐の個数を二つ増やす。
操作の回数はそれぞれ $10,8,9$ なので、最終扇の最大錐は全て24個である。
これらは24個の大小関係の錐によって細分されるため、各最大錐が一つずつ一致する。
同じ結論は、最終最大錐が
\[
 \bigl\langle
 v_{i_1},\ v_{i_1}+v_{i_2},\ v_{i_1}+v_{i_2}+v_{i_3}
 \bigr\rangle,
 \qquad (i_1,i_2,i_3,i_4)\in\mathfrak S_4
\]
となることを上の置換規則で列挙しても得られる。
付属の \texttt{check\_birational.py} は中心の存在、全段階の滑らかさと面の条件、
最終24錐の集合の一致を整数演算で検証する。
Poisson 構造については、各中間模型が同じトーラスの滑らかなコンパクト化なので、
三つのトーラス作用のベクトル場の外積がそのまま正則な延長を与える。
\end{proof}

\subsection{\texorpdfstring{$(\mathbb P^1)^3$}{P1 x P1 x P1} の対数型からトーラス型への変換}
$S=\mathbb P^1_x\times\mathbb P^1_z$ とし、
$F,G\in H^0(S,\mathcal O_S(1,1))$ を一般の切断とする。
第三因子の座標を $y$ とし、対応する対数型の双ベクトルを局所体積形式
$dx\wedge dy\wedge dz$ との縮約で
\[
 \alpha=yFG\left(a\frac{dy}{y}+b\,d\log\frac FG\right)
\]
と表す。これは $y=0,\infty,F=0,G=0$ に対応する対数型である。
各係数は
\[
 \alpha_x=by(GF_x-FG_x),\quad
 \alpha_y=aFG,\quad
 \alpha_z=by(GF_z-FG_z)
\]
であり、それぞれの多重次数は高々 $(0,1,2)$、$(2,0,2)$、$(2,1,0)$
となる。従って対応する双ベクトルは三重積全体で正則である。
$F=xA(z)+B(z)$、$G=xC(z)+D(z)$ と書き、
$\Delta=AD-BC\in H^0(\mathbb P^1_z,\mathcal O(2))$ と置く。
有理写像
\[
 \Phi(x,y,z)=\left(u=\frac FG,\,y,\,z\right)
\]
は、$\Delta\ne0$ の開集合上で
$x=(B-uD)/(uC-A)$ を逆写像に持つ。
$\partial u/\partial x=\Delta/G^2$ であるから、体積形式の変換も含めて
\[
 \iota_{\Phi_*\pi}(du\wedge dy\wedge dz)
   =\frac{\Delta}{G^2}\alpha
   =\Delta\bigl(a u\,dy+b y\,du\bigr).
\]
従って双ベクトルそのものの像は
\[
 \Phi_*\pi
   =\Delta(z)\left(b y\partial_y\wedge\partial_z
                         +a u\partial_z\wedge\partial_u\right).
 \tag{T}\label{eq:birational-log-toric}
\]
直接に Poisson 括弧を計算しても
\[
 \{u,y\}=0,\qquad \{y,z\}=by\Delta,\qquad
 \{z,u\}=au\Delta
\]
を得る。像の各項も三つの因子上の大域正則ベクトル場の外積である。
一般の $F,G$ に対して $\Delta\partial_z$ は相異なる二つの零点を持ち、
適当な射影座標 $w$ で $\Delta\partial_z=\kappa w\partial_w$ と書ける。
式 \eqref{eq:birational-log-toric} はこの座標でトーラス型となる。
その係数行列は $\lambda_{uy}=0$ を満たすので、この構成の像は
トーラス型のうち一つの係数が零となる部分族である。

この双有理写像も Poisson blow-up の組合せで実現できる。
一般の $F,G$ の共通零点は $S$ 上の相異なる二点 $p_1,p_2$ であり、
その $z$ 座標は $\Delta$ の二つの零点である。
$X=S\times\mathbb P^1_y$ 上の不定点集合は
\[
 E=(F=G=0)=\bigl(\{p_1\}\sqcup\{p_2\}\bigr)\times\mathbb P^1_y.
\]
$E$ は滑らかな余次元二の零 Poisson 部分多様体である。
グラフの閉包は $\mathrm{Bl}_E X$ で、方程式は
$u_0F-u_1G=0$ である。反対側から見ると、これは
\[
 E'=\{u_0A-u_1C=u_0B-u_1D=0\}
       \subset\mathbb P^1_u\times\mathbb P^1_z\times\mathbb P^1_y
\]
の blow-up である。$E'$ も二本の相異なる $\mathbb P^1_y$ の和であり、
式 \eqref{eq:birational-log-toric} はその上で零となる。
滑らかな余次元二の Poisson 中心では conormal Lie algebra の退化条件は
自動的に成り立つので、Polishchuk の blow-up 定理によって両側の
blow-up に正則 Poisson 構造が持ち上がる。両構造は共通の稠密開集合で
一致するため、グラフ全体で一致する。これにより、両コンパクト化上の正則な双ベクトルと共通 Poisson 解消が得られる。


\subsection{\texorpdfstring{$X_{12}$}{X12} の有理 pencil 成分}
分類節の $C_{(1,1)^2(0,1)}$ を考える。
$F,G\in H^0(X_{12},\mathcal O(1,1))$、
$L\in H^0(\PP^2,\mathcal O(1))$ に対し
\[
 \alpha=L(F\,dG-G\,dF)
\]
と置く。$F=xa(v,w)+b(v,w)$、$G=xc(v,w)+d(v,w)$ と書くと、
$u=G/F$ は Casimir であり
\[
 \{v,w\}=L(bc-ad),\qquad \{u,v\}=\{u,w\}=0.
\]
従って $(x,v,w)\mapsto(u,v,w)$ は、零三次曲線が直線と二次曲線の和である
垂直型への Poisson 双有理写像である。
一般の係数では、行列式 $bc-ad=0$ は滑らかな二次曲線であり、
$a,b,c,d$ は平面上で同時には零にならない。
$F=G=0$ および逆写像の基底曲線は、この二次曲線と同型な滑らかな有理曲線である。
両曲線上でそれぞれの双ベクトルは零になる。
前述の比の閉グラフの議論と命題\ref{prop:poisson-centres}により、
そのグラフは両側の Poisson blow-up である。

より一般に $\CC(S)(x)$ の括弧が $\CC(S)$ 上の括弧 $\sigma$ を保ち、
$u=(ax+b)/(cx+d)$ が非定数の Casimir ならば、
$(x,s)\mapsto(u,s)$ は同じ $\sigma$ を垂直に持つ積への Poisson 双有理写像を与える。
正則な共通模型の構成では、基底イデアルと各中心上の双ベクトルを別に確認する。

\subsection{共通模型上の古典的 SOD}

滑らかな中心 $Y\subset X$ の余次元を $c$、$f:\widetilde X\to X$ をその blow-up、
$j:E\hookrightarrow\widetilde X$、$p:E\to Y$ を例外因子の射とする。
Orlov \cite[定理4.3]{OrlovBlowup} の式は
\[
 \operatorname{Perf}(\widetilde X)=
 \left\langle
 j_*\bigl(p^*\operatorname{Perf}(Y)\otimes\mathcal O_E(1-c)\bigr),\ldots,
 j_*\bigl(p^*\operatorname{Perf}(Y)\otimes\mathcal O_E(-1)\bigr),
 Lf^*\operatorname{Perf}(X)
 \right\rangle.
\]
三次元では曲線中心につき $\operatorname{Perf}(Y)$ が一個、
点中心につき $\operatorname{Perf}(\CC)$ が二個現れる。
例えば定理\ref{thm:rat22-vertical}の共通模型では、埋込み関手を省略すると
\[
\begin{aligned}
 \operatorname{Perf}(Z)
 &=\left\langle\operatorname{Perf}(E'),\operatorname{Perf}(\CC),
                  \operatorname{Perf}(\CC),\operatorname{Perf}(\PP^3)\right\rangle,\\
 \operatorname{Perf}(Z)
 &=\left\langle\operatorname{Perf}(\Gamma),
                  \operatorname{Perf}(\PP^1\times\PP^2)\right\rangle.
\end{aligned}
\]
$E'\cong\Gamma$ は楕円曲線である。
三重積 pencil の共通模型でも、両側にそれぞれ楕円曲線の圏が一個加わる。

命題\ref{prop:toric-explicit-blowups}の $Z_{\rm per}$ は長さ24の full exceptional collection を持つ。
実際、三つの端点の標準的な積の例外集の長さは順に $4,6,8$ である。
各点 blow-up には二つの例外対象が加わり、各曲線中心は滑らかな完備トーリック曲線
$\PP^1$ なので、その二項例外集が加わる。従って三つの構成から得られる長さは
\[
 4+2\cdot4+2\cdot6
 =6+2\cdot1+2\cdot8
 =8+2\cdot2+2\cdot6=24
\]
と一致する。この記述は上の可換な共通模型の導来圏についてのものである。
これらの式の非可換化には、中心とその例外部分を含む blow-up 圏の構成が必要になる。
その構成に適用できる先行結果を第\ref{sec:nc-blowups}節で整理する。
